\chapter{VBE formulių sąrašas}

Šiame skyriuje pateikiamos visos formulės, kuriomis galima naudotis matematikos
valstybinio brandos egzamino metu. Formulės atitinka Nacionalinės švietimo agentūros
(NŠA) parengtą oficialų formulių rinkinį.

% =====================================================================
\section{Greitosios daugybos formulės}
% =====================================================================

\begin{tcolorbox}[colback=blue!5!white, colframe=blue!60!black, title=Greitosios daugybos formulės]
\begin{align}
(a \pm b)^2 &= a^2 \pm 2ab + b^2 \\[6pt]
(a \pm b)^3 &= a^3 \pm 3a^2b + 3ab^2 \pm b^3 \\[6pt]
a^2 - b^2   &= (a-b)(a+b) \\[6pt]
a^3 \pm b^3 &= (a \pm b)(a^2 \mp ab + b^2)
\end{align}
\end{tcolorbox}

% =====================================================================
\section{Laipsniai ir šaknys}
% =====================================================================

\subsection{Laipsnių savybės}

Čia $a > 0$, $a \neq 1$, $n \in \mathbb{N}$, $n > 1$, $m \in \mathbb{Z}$.

\begin{align}
a^{-n}        &= \frac{1}{a^n} \\[6pt]
a^m \cdot a^n &= a^{m+n} \\[6pt]
a^m : a^n     &= a^{m-n} \\[6pt]
(a^m)^n       &= a^{mn} \\[6pt]
(ab)^n        &= a^n b^n \\[6pt]
\left(\frac{a}{b}\right)^n &= \frac{a^n}{b^n}
\end{align}

\subsection{Šaknų savybės}

Jei $\sqrt[n]{b} = a$, tai $a^n = b$ \quad ($n \in \mathbb{N}$, $n > 1$):

\begin{align}
\sqrt[n]{ab}         &= \sqrt[n]{a}\cdot\sqrt[n]{b} \\[6pt]
\sqrt[n]{\frac{a}{b}} &= \frac{\sqrt[n]{a}}{\sqrt[n]{b}} \\[6pt]
\sqrt[n]{a^m}        &= a^{m/n} \\[6pt]
\sqrt[n]{\sqrt[k]{a}} &= \sqrt[nk]{a}
\end{align}

% =====================================================================
\section{Logaritmai}
% =====================================================================

Jei $\log_a b = c$, tai $a^c = b$ \quad ($a > 0$, $a \neq 1$, $b > 0$):

\begin{tcolorbox}[colback=green!5!white, colframe=green!60!black, title=Logaritmų savybės]
\begin{align}
\log_a(bc)  &= \log_a b + \log_a c \\[6pt]
\log_a\!\left(\frac{b}{c}\right) &= \log_a b - \log_a c \\[6pt]
\log_a b^k  &= k\log_a b \\[6pt]
\log_b c    &= \frac{\log_a c}{\log_a b} \quad \text{(perėjimo formulė)}
\end{align}
\end{tcolorbox}

\begin{remark}
Natūrinis logaritmas: $\ln x = \log_e x$, kur $e \approx 2{,}718$. \\
Dešimtainis logaritmas: $\lg x = \log_{10} x$.
\end{remark}

% =====================================================================
\section{Trigonometrija}
% =====================================================================

\subsection{Pagrindiniai sąryšiai}

\begin{align}
\sin^2\alpha + \cos^2\alpha &= 1 \\[6pt]
\tg\alpha   &= \frac{\sin\alpha}{\cos\alpha} \\[6pt]
\ctg\alpha  &= \frac{\cos\alpha}{\sin\alpha} \\[6pt]
1 + \tg^2\alpha  &= \frac{1}{\cos^2\alpha} \\[6pt]
1 + \ctg^2\alpha &= \frac{1}{\sin^2\alpha}
\end{align}

\subsection{Dvigubo kampo formulės}

\begin{align}
\sin 2\alpha &= 2\sin\alpha\cos\alpha \\[6pt]
\cos 2\alpha &= \cos^2\alpha - \sin^2\alpha \\[6pt]
\tg 2\alpha  &= \frac{2\tg\alpha}{1-\tg^2\alpha}
\end{align}

\subsection{Pusės kampo formulės}

\begin{align}
2\sin^2\alpha &= 1 - \cos 2\alpha \\[6pt]
2\cos^2\alpha &= 1 + \cos 2\alpha
\end{align}

\subsection{Sumos ir skirtumo formulės}

\begin{align}
\sin(\alpha \pm \beta) &= \sin\alpha\cos\beta \pm \cos\alpha\sin\beta \\[6pt]
\cos(\alpha \pm \beta) &= \cos\alpha\cos\beta \mp \sin\alpha\sin\beta \\[6pt]
\tg(\alpha \pm \beta)  &= \frac{\tg\alpha \pm \tg\beta}{1 \mp \tg\alpha\tg\beta}
\end{align}

\subsection{Trigonometrinių funkcijų reikšmių lentelė}

\begin{center}
\begin{tabular}{|c|c|c|c|c|c|}
\hline
$\alpha$ (laipsniais) & $0°$ & $30°$ & $45°$ & $60°$ & $90°$ \\
\hline
$\alpha$ (radianais)  & $0$ & $\dfrac{\pi}{6}$ & $\dfrac{\pi}{4}$ & $\dfrac{\pi}{3}$ & $\dfrac{\pi}{2}$ \\[8pt]
\hline
$\sin\alpha$ & $0$ & $\dfrac{1}{2}$ & $\dfrac{\sqrt{2}}{2}$ & $\dfrac{\sqrt{3}}{2}$ & $1$ \\[8pt]
\hline
$\cos\alpha$ & $1$ & $\dfrac{\sqrt{3}}{2}$ & $\dfrac{\sqrt{2}}{2}$ & $\dfrac{1}{2}$ & $0$ \\[8pt]
\hline
$\tg\alpha$  & $0$ & $\dfrac{\sqrt{3}}{3}$ & $1$ & $\sqrt{3}$ & $-$ \\[8pt]
\hline
\end{tabular}
\end{center}

\subsection{Trigonometrinės lygtys}

\begin{tcolorbox}[colback=orange!5!white, colframe=orange!60!black, title=Trigonometrinių lygčių sprendiniai]
\begin{align}
\sin x = a \quad &\Rightarrow \quad x = (-1)^k \arcsin a + \pi k, \quad k \in \mathbb{Z},
    \quad a \in [-1;\,1] \\[6pt]
\cos x = a \quad &\Rightarrow \quad x = \pm\arccos a + 2\pi k, \quad k \in \mathbb{Z},
    \quad a \in [-1;\,1] \\[6pt]
\tg x = a  \quad &\Rightarrow \quad x = \arctg a + \pi k, \quad k \in \mathbb{Z},
    \quad a \in \mathbb{R}
\end{align}
\end{tcolorbox}

% =====================================================================
\section{Sekos ir procentai}
% =====================================================================

\subsection{Aritmetinė progresija}

Čia $a_1$ – pirmasis narys, $d$ – skirtumas, $n$ – nario eilės numeris,
$S_n$ – pirmųjų $n$ narių suma.

\begin{align}
a_n &= a_1 + d(n-1) \\[6pt]
d   &= a_{n+1} - a_n \\[6pt]
S_n &= \frac{a_1 + a_n}{2} \cdot n = \frac{2a_1 + d(n-1)}{2} \cdot n
\end{align}

\subsection{Geometrinė progresija}

Čia $b_1 \neq 0$ – pirmasis narys, $q \neq 0$ – vardiklis, $S_n$ – pirmųjų $n$ narių
suma, $S$ – nykstamosios geometrinės progresijos suma.

\begin{align}
b_n &= b_1 \cdot q^{n-1} \\[6pt]
q   &= \frac{b_{n+1}}{b_n} \\[6pt]
S_n &= \frac{b_1(1 - q^n)}{1-q} = \frac{b_1 - q b_n}{1-q}, \quad q \neq 1 \\[6pt]
S   &= \frac{b_1}{1-q}, \quad |q| < 1 \quad \text{(nykstamoji progresija)}
\end{align}

\subsection{Sudėtiniai procentai}

\begin{tcolorbox}[colback=purple!5!white, colframe=purple!60!black, title=Sudėtinių procentų formulė]
\[
S_n = S_0 \left(1 + \frac{p}{100}\right)^n
\]
\end{tcolorbox}

čia $S_0$ – pradinė reikšmė, $p$ – procentų skaičius, $n$ – laikotarpių skaičius.

% =====================================================================
\section{Planimetrinės figūros}
% =====================================================================

\subsection{Trikampis}

Čia $a$, $b$, $c$ – kraštinių ilgiai; $A$, $B$, $C$ – prieš jas esančių kampų
didumai; $p = \dfrac{a+b+c}{2}$ – pusperimetris; $r$ – įbrėžtinio,
$R$ – apibrėžtinio apskritimo spindulys; $h_a$ – aukštinė į kraštinę $a$.

\textbf{Kosinusų teorema:}
\[
a^2 = b^2 + c^2 - 2bc\cos A
\]

\textbf{Sinusų teorema:}
\[
\frac{a}{\sin A} = \frac{b}{\sin B} = \frac{c}{\sin C} = 2R
\]

\textbf{Trikampio plotas:}
\[
S = \frac{1}{2}ah_a = \frac{1}{2}ab\sin C
  = \sqrt{p(p-a)(p-b)(p-c)}
  = rp
  = \frac{abc}{4R}
\]

\subsection{Skritulio išpjova}

Čia $R$ – spindulio ilgis, $\alpha$ – centrinio kampo didumas laipsniais.

\begin{align}
S_{\text{išpj.}} &= \frac{\alpha}{360°}\pi R^2 \\[6pt]
C_{\text{išpj.}} &= \frac{\alpha}{360°} \cdot 2\pi R
\end{align}

% =====================================================================
\section{Erdvinės figūros}
% =====================================================================

\subsection{Ritinys}

Čia $R$ – pagrindo spindulio ilgis, $H$ – aukštinės ilgis.
\begin{align}
S_{\text{šon.}} &= 2\pi R H \\[6pt]
V &= \pi R^2 H
\end{align}

\subsection{Kūgis}

Čia $R$ – pagrindo spindulio ilgis, $l$ – sudaromosios ilgis, $H$ – aukštinės ilgis.
\begin{align}
S_{\text{šon.}} &= \pi R l \\[6pt]
V &= \frac{1}{3}\pi R^2 H
\end{align}

\subsection{Nupjautinis kūgis}

Čia $R$ ir $r$ – pagrindų spindulių ilgiai, $l$ – sudaromosios ilgis, $H$ – aukštinės ilgis.
\begin{align}
S_{\text{šon.}} &= \pi(R + r)l \\[6pt]
V &= \frac{\pi H}{3}(R^2 + Rr + r^2)
\end{align}

\subsection{Rutulys}

Čia $R$ – spindulio ilgis.
\begin{align}
S &= 4\pi R^2 \\[6pt]
V &= \frac{4}{3}\pi R^3
\end{align}

\subsection{Rutulio nuopjova}

Čia $R$ – rutulio spindulio ilgis, $H$ – nuopjovos aukštis.
\begin{align}
S &= 2\pi R H \\[6pt]
V &= \frac{\pi H^2(3R - H)}{3}
\end{align}

\subsection{Piramidė}

Čia $S_{\text{pag.}}$ – pagrindo plotas, $H$ – aukštinės ilgis.
\[
V = \frac{1}{3} S_{\text{pag.}} \cdot H
\]

\subsection{Nupjautinė piramidė}

Čia $S_1$, $S_2$ – pagrindų plotai, $H$ – aukštinės ilgis.
\[
V = \frac{H}{3}\!\left(S_1 + \sqrt{S_1 S_2} + S_2\right)
\]

% =====================================================================
\section{Vektoriai}
% =====================================================================

\subsection{Vektoriaus ilgis erdvėje}

Čia $\vec{a} = (x;\, y;\, z)$:
\[
|\vec{a}| = \sqrt{x^2 + y^2 + z^2}
\]

\subsection{Skaliarinė sandauga}

Čia $\vec{a} = (x_1;\, y_1;\, z_1)$, $\vec{b} = (x_2;\, y_2;\, z_2)$,
$\alpha$ – kampas tarp vektorių:
\[
\vec{a} \cdot \vec{b} = x_1 x_2 + y_1 y_2 + z_1 z_2 = |\vec{a}|\,|\vec{b}|\cos\alpha
\]

% =====================================================================
\section{Išvestinės}
% =====================================================================

\subsection{Diferencijavimo taisyklės}

Čia $u$ ir $v$ – diferencijuojamosios funkcijos, $c$ – konstanta.

\begin{tcolorbox}[colback=red!5!white, colframe=red!60!black, title=Diferencijavimo taisyklės]
\begin{align}
(cu)'           &= cu' \\[4pt]
(u \pm v)'      &= u' \pm v' \\[4pt]
(uv)'           &= u'v + uv' \\[4pt]
\left(\frac{u}{v}\right)' &= \frac{u'v - uv'}{v^2} \\[6pt]
(x^n)'          &= nx^{n-1}
\end{align}
\end{tcolorbox}

\subsection{Pagrindinių funkcijų išvestinės}

\begin{center}
\begin{tabular}{|c|c|}
\hline
\textbf{Funkcija} $f(x)$ & \textbf{Išvestinė} $f'(x)$ \\
\hline
$x^n$          & $nx^{n-1}$ \\[4pt]
\hline
$\sqrt{x}$     & $\dfrac{1}{2\sqrt{x}}$ \\[8pt]
\hline
$a^x$          & $a^x \ln a$ \\[4pt]
\hline
$e^x$          & $e^x$ \\[4pt]
\hline
$\log_a x$     & $\dfrac{1}{x \ln a}$ \\[8pt]
\hline
$\ln x$        & $\dfrac{1}{x}$ \\[8pt]
\hline
$\sin x$       & $\cos x$ \\[4pt]
\hline
$\cos x$       & $-\sin x$ \\[4pt]
\hline
$\tg x$        & $\dfrac{1}{\cos^2 x}$ \\[8pt]
\hline
$\ctg x$       & $-\dfrac{1}{\sin^2 x}$ \\[8pt]
\hline
\end{tabular}
\end{center}

\subsection{Sudėtinės funkcijos išvestinė}

Jei $h(x) = g(f(x))$, tai:
\[
h'(x) = g'(f(x)) \cdot f'(x)
\]

\subsection{Liestinės lygtis}

Funkcijos $y = f(x)$ grafiko liestinės, liečiančios grafiką taške
$\bigl(x_0;\, f(x_0)\bigr)$, lygtis:
\[
y = f(x_0) + f'(x_0)(x - x_0)
\]
čia $f'(x_0)$ – liestinės krypties koeficientas.

% =====================================================================
\section{Kombinatorika}
% =====================================================================

\begin{align}
\text{Derinių skaičius:} \quad
C_n^k &= C_n^{n-k} = \frac{n!}{k!\,(n-k)!} \\[10pt]
\text{Gretinių skaičius:} \quad
A_n^k &= \frac{n!}{(n-k)!}
\end{align}

\begin{remark}
Pagal susitarimą $0! = 1$.
\end{remark}

% =====================================================================
\section{Tikimybių teorija ir statistika}
% =====================================================================

\subsection{Atsitiktinio dydžio charakteristikos}

Tegu atsitiktinis dydis $X$ įgyja reikšmes $x_1, x_2, \ldots, x_n$
su atitinkamomis tikimybėmis $p_1, p_2, \ldots, p_n$.

\textbf{Matematinė viltis:}
\[
EX = x_1 p_1 + x_2 p_2 + \cdots + x_n p_n
\]

\textbf{Dispersija:}
\[
DX = (x_1 - EX)^2 p_1 + (x_2 - EX)^2 p_2 + \cdots + (x_n - EX)^2 p_n
\]

\textbf{Standartinis nuokrypis:}
\[
\sigma = \sqrt{DX}
\]

\subsection{Bernulio formulė}

Jei eksperimentas kartojamas $n$ kartų nepriklausomai, sėkmės tikimybė kiekviename
bandyme lygi $p$ (nesėkmės $q = 1-p$), tai tikimybė gauti lygiai $k$ sėkmių:
\[
P(X = k) = C_n^k \, p^k \, q^{n-k}, \quad k = 0, 1, \ldots, n
\]

